\chapter{The decoder, pixel shuffle and one full forward trace}\label{ch:decoder}

\begin{onething}
The decoder brings the four encoder maps to one width (256 channels, by 1 $\times$ 1 convolutions) and
one size (32 $\times$ 32, by bilinear interpolation), stacks and fuses them, refines at quarter
resolution, and then reaches full resolution by \emph{rearranging} 512 channels into 4 $\times$ 4 blocks
of 32 channels, so that every output pixel gets its own learned prediction.
\end{onething}

\Cref{ch:encoder} ended with four feature maps of the same ground: F1 \shape{(B, 32, 32, 32)}, F2
\shape{(B, 64, 16, 16)}, F3 \shape{(B, 160, 8, 8)} and F4 \shape{(B, 256, 4, 4)}. The model must return a
reconstruction of the input's shape, \shape{(B, 165, 128, 128)}, because a residual is a subtraction and
you can only subtract tensors of the same shape (\cref{ch:recon}). This chapter builds that last stretch:
the decoder (\code{app/foundation_models/components/seg_former_decoder.py}, 176 lines), the
decompressor and the denormalisation. It ends by walking one batch through the whole model.

\section{The problem, and why the obvious fixes fail}\label{sec:decoder:why}

Three things must happen: the four maps have to be combined, the result has to reach 128 $\times$ 128,
and the 32 compressed channels have to become 165 bands again. The obvious ways to do the first two
each fail.

\begin{itemize}
\item \textbf{Use only F4 and upsample it 32 times.} Sixteen tokens cannot say where a single odd pixel
is; all fine detail is lost.
\item \textbf{Run convolutions at full resolution} (a U-Net style decoder). It works, but a 3 $\times$ 3
convolution at 128 $\times$ 128 costs sixteen times as much as at 32 $\times$ 32, because there are
$128^2 / 32^2 = 16$ times as many positions.
\item \textbf{Work at 32 $\times$ 32 and finish with bilinear upsampling.} Cheap, but bilinear
interpolation \emph{averages} neighbours, so a single anomalous pixel's prediction is blended with its
fifteen ordinary neighbours, and the residual that should flag it is diluted. That is exactly the
signal this system measures.
\end{itemize}

The SegFormer decoder \citep{xie2021segformer}, as modified here, takes a fourth route: do all the
expensive thinking at 32 $\times$ 32, and jump to full size with a free rearrangement,
\index{pixel shuffle}\emph{pixel shuffle} \citep{shi2016realtime}, that gives each output pixel its own
channel and therefore its own weights.

\begin{analogy}
A mixing desk combines a 48-track recording into one stereo master. Each track first goes through its
own fader and equaliser so that all tracks sit at a common level (the 1 $\times$ 1 projections to 256
channels). Tracks recorded at different sample rates are resampled to one rate (the bilinear
upsampling). The desk sums them into a mix bus (the concatenation and the fuse convolution) and a
mastering stage polishes the result (the refine convolutions). The last step is the one the analogy
makes vivid: a surround-sound master is not ``stretched'' to more speakers; each speaker gets its own
dedicated channel. Pixel shuffle is that: 16 channels per position, one for each of the 16 output
pixels in the 4 $\times$ 4 block.
\end{analogy}

\section{Step 1: one width for all four maps}\label{sec:decoder:project}

You cannot stack maps of 32, 64, 160 and 256 channels into one tensor and treat every channel alike, so
each map first passes through its own \index{convolution!1x1}1 $\times$ 1 convolution to the common
\index{decoder width}\emph{decoder width} $D_{\text{dec}} = 256$:
\begin{lstlisting}
self.linear_projections = nn.ModuleList([
    nn.Conv2d(embed_dim, decoder_dim, kernel_size=1)    # 32/64/160/256 -> 256
    for embed_dim in embed_dims])
\end{lstlisting}
A 1 $\times$ 1 convolution looks at one position at a time and mixes only channels: at every position
$(h, w)$ it computes
\begin{equation}\label{eq:decoder:proj}
\vect{p}_{h,w} = \mat{W}_i\, \vect{f}_{h,w} + \vect{b}_i, \qquad \mat{W}_i \in \R^{256 \times C_i},\ \vect{b}_i \in \R^{256},
\end{equation}
with the same $\mat{W}_i$ and $\vect{b}_i$ at every position (\cref{ch:shapes}). It is a linear layer
applied per pixel.

\begin{worked}[label={we:decoder:proj}]{A 1 \texorpdfstring{$\times$}{x} 1 projection at one position}
Shrink the sizes: 2 input channels, 3 output channels. At one position the feature vector is
$\vect{f} = (2, -1)\T$. Let
\[
\mat{W} = \begin{pmatrix} 1 & 0\\ 0.5 & 2\\ -1 & 1 \end{pmatrix}, \qquad \vect{b} = \begin{pmatrix} 0\\ 1\\ 0.5 \end{pmatrix}
\]
(illustrative weights). Multiply row by row, then add the bias. Do it on paper:
\begin{align*}
p_0 &= 1 \cdot 2 + 0 \cdot (-1) + 0 = 2 + 0 + 0 = 2,\\
p_1 &= 0.5 \cdot 2 + 2 \cdot (-1) + 1 = 1 - 2 + 1 = 0,\\
p_2 &= -1 \cdot 2 + 1 \cdot (-1) + 0.5 = -2 - 1 + 0.5 = -2.5 .
\end{align*}
So $\vect{p} = (2, 0, -2.5)\T$. The decoder does the same with $\mat{W}_1 \in \R^{256 \times 32}$ at each
of F1's $32 \times 32 = 1024$ positions, and similarly for F2--F4.
\end{worked}

Shapes after step 1:
\[
\begin{array}{lll}
\text{F1}: (B, 32, 32, 32) & \to & (B, 256, 32, 32)\\
\text{F2}: (B, 64, 16, 16) & \to & (B, 256, 16, 16)\\
\text{F3}: (B, 160, 8, 8) & \to & (B, 256, 8, 8)\\
\text{F4}: (B, 256, 4, 4) & \to & (B, 256, 4, 4)
\end{array}
\]
F4 already has 256 channels but still gets its own projection: the four maps live in different learned
spaces and each needs its own mapping into the shared one.

\section{Step 2: one size for all four maps, by bilinear interpolation}\label{sec:decoder:bilinear}

Every projected map is resized to F1's 32 $\times$ 32:
\begin{lstlisting}
if proj.shape[2] != target_h or proj.shape[3] != target_w:          # F1 is skipped
    proj = F.interpolate(proj, size=(target_h, target_w),
                         mode='bilinear', align_corners=False)      # (B, 256, 32, 32)
\end{lstlisting}
F2 grows by a factor 2, F3 by 4, F4 by 8. What does \emph{bilinear} compute?

\subsection{Linear interpolation on a line}

\begin{definition}[Linear interpolation]
Given values $f_0$ at position 0 and $f_1$ at position 1, the \index{interpolation!linear}\emph{linear
interpolation} at a position $u \in [0, 1]$ between them is
\begin{equation}\label{eq:decoder:lerp}
f(u) = (1 - u)\, f_0 + u\, f_1 .
\end{equation}
\end{definition}
Derive it: the straight line through $(0, f_0)$ and $(1, f_1)$ has slope $(f_1 - f_0)/(1 - 0) = f_1 -
f_0$ and starts at $f_0$, so $f(u) = f_0 + u (f_1 - f_0)$. Expand the bracket: $f_0 + u f_1 - u f_0$.
Collect the $f_0$ terms: $(1 - u) f_0 + u f_1$. The two weights $1 - u$ and $u$ are non-negative and sum
to 1, so $f(u)$ is a weighted average of the two neighbours, weighted by closeness. At $u = 0$ you get
$f_0$, at $u = 1$ you get $f_1$, at $u = 0.25$ you get three parts $f_0$ to one part $f_1$.

\subsection{Where does an output pixel sit on the input grid?}

To resize a row of $n_{\text{in}}$ values to $n_{\text{out}}$ values you must decide where each output
pixel falls among the inputs. With \code{align_corners=False}, PyTorch treats every pixel as a cell of
width 1 whose \emph{centre} is the sample, and lines up the outer edges of the two grids. Output pixel
$d$ has its centre at $d + \frac12$ in output units; multiplying by the scale $s = n_{\text{in}} /
n_{\text{out}}$ converts to input units; subtracting $\frac12$ converts from ``centre'' back to an index:
\begin{equation}\label{eq:decoder:src}
u_d = \Big(d + \tfrac12\Big)\, \frac{n_{\text{in}}}{n_{\text{out}}} - \tfrac12, \qquad \text{clamped below at } 0 .
\end{equation}
The output value is then the linear interpolation between input indices $\lfloor u_d \rfloor$ and
$\lfloor u_d \rfloor + 1$ with weight $u_d - \lfloor u_d \rfloor$; when $\lfloor u_d \rfloor + 1$ falls
off the end, the last input value is used.

\begin{worked}[label={we:decoder:lerp}]{Upsampling a row of two values to four}
Input $[0, 4]$, so $f_0 = 0$, $f_1 = 4$; $n_{\text{in}} = 2$, $n_{\text{out}} = 4$, scale $s = 2/4 =
0.5$. On paper, one output at a time:
\begin{align*}
d = 0:&\quad u = (0.5)(0.5) - 0.5 = -0.25 \to \text{clamp to } 0, & \text{value} &= f_0 = 0\\
d = 1:&\quad u = (1.5)(0.5) - 0.5 = 0.25, & \text{value} &= 0.75 \cdot 0 + 0.25 \cdot 4 = 1\\
d = 2:&\quad u = (2.5)(0.5) - 0.5 = 0.75, & \text{value} &= 0.25 \cdot 0 + 0.75 \cdot 4 = 3\\
d = 3:&\quad u = (3.5)(0.5) - 0.5 = 1.25, & \text{value} &= f_1 = 4
\end{align*}
(At $d = 3$ there is no input index 2 to interpolate towards, so the last value is used.)
Output $[0, 1, 3, 4]$. \code{F.interpolate} returns exactly this
(\file{research/textbook/tools/compute_masking_output.txt}).
\end{worked}

\subsection{Bilinear: two linear interpolations}

On a grid, a point at fractional position $(u, v)$ (row, column) sits inside a square of four samples
$f_{00}, f_{01}$ (upper row) and $f_{10}, f_{11}$ (lower row), with fractional offsets $\alpha = u -
\lfloor u \rfloor$ down and $\beta = v - \lfloor v \rfloor$ across. \index{interpolation!bilinear}\emph{Bilinear}
interpolation interpolates along each row first, then between the two results:
\begin{align}
\text{upper row:}\quad g_0 &= (1 - \beta) f_{00} + \beta f_{01},\notag\\
\text{lower row:}\quad g_1 &= (1 - \beta) f_{10} + \beta f_{11},\notag\\
\text{between rows:}\quad f &= (1 - \alpha)\, g_0 + \alpha\, g_1 .\label{eq:decoder:bilinear}
\end{align}
Substitute $g_0$ and $g_1$ and expand to see all four weights at once:
\[
f = (1-\alpha)(1-\beta) f_{00} + (1-\alpha)\beta f_{01} + \alpha(1-\beta) f_{10} + \alpha\beta f_{11} .
\]
The four weights are products of the two one-dimensional weights and sum to $\big((1-\alpha)+\alpha\big)
\big((1-\beta)+\beta\big) = 1$. The expansion is symmetric in the two directions, so interpolating the
columns first gives the same answer.

\begin{worked}[label={we:decoder:bilinear}]{Bilinear upsampling of a 2 \texorpdfstring{$\times$}{x} 2 grid to 4 \texorpdfstring{$\times$}{x} 4}
Input $\begin{pmatrix} 1 & 3\\ 5 & 7\end{pmatrix}$. Rows and columns each go from 2 to 4, so by
\cref{we:decoder:lerp} the fractional positions are $0, 0.25, 0.75, 1$ (after clamping) in both
directions. First interpolate along each input row, using the weights of \cref{we:decoder:lerp}:
\begin{align*}
\text{row } [1, 3]:&\quad [\,1,\ 0.75 \cdot 1 + 0.25 \cdot 3,\ 0.25 \cdot 1 + 0.75 \cdot 3,\ 3\,] = [1,\ 1.5,\ 2.5,\ 3],\\
\text{row } [5, 7]:&\quad [\,5,\ 0.75 \cdot 5 + 0.25 \cdot 7,\ 0.25 \cdot 5 + 0.75 \cdot 7,\ 7\,] = [5,\ 5.5,\ 6.5,\ 7].
\end{align*}
Now interpolate down each column between these two rows, with the same weights $0, 0.25, 0.75, 1$. For
the second output column (values $1.5$ above, $5.5$ below):
\begin{align*}
\text{out row 0}&: 1.5, \\
\text{out row 1}&: 0.75 \cdot 1.5 + 0.25 \cdot 5.5 = 1.125 + 1.375 = 2.5,\\
\text{out row 2}&: 0.25 \cdot 1.5 + 0.75 \cdot 5.5 = 0.375 + 4.125 = 4.5,\\
\text{out row 3}&: 5.5 .
\end{align*}
Doing every column the same way:
\[
\begin{pmatrix} 1 & 1.5 & 2.5 & 3\\ 2 & 2.5 & 3.5 & 4\\ 4 & 4.5 & 5.5 & 6\\ 5 & 5.5 & 6.5 & 7 \end{pmatrix}.
\]
This is exactly what \code{F.interpolate(..., mode='bilinear', align_corners=False)} returns (checked in
\file{research/textbook/tools/compute_masking.py}). Notice that no output value is outside the range
$[1, 7]$ of the inputs, and that no output differs from its neighbour by more than 2: bilinear
interpolation can only smooth.
\end{worked}

That last remark is why bilinear interpolation is fine \emph{here}, on smooth coarse context from F2--F4,
and wrong as the final step to full resolution (\cref{sec:decoder:shuffle}).

\section{Steps 3 and 4: stack, then fuse}\label{sec:decoder:fuse}

\begin{lstlisting}
fused = torch.cat(unified, dim=1)     # 4 x (B, 256, 32, 32) -> (B, 1024, 32, 32)
fused = self.fuse_conv(fused)         # Conv2d(1024, 256, k=1) -> (B, 256, 32, 32)
fused = self.fuse_act(fused)          # GELU
\end{lstlisting}
\index{concatenation}\index{fuse convolution}\code{torch.cat(..., dim=1)} glues along the channel axis: $4 \times 256 = 1024$ channels. At every
one of the 1024 positions there are now 1024 numbers: 256 describing fine texture there (from F1), 256
medium structure (F2), 256 coarse (F3) and 256 global context (F4). The fuse convolution is another 1
$\times$ 1 convolution, \eqref{eq:decoder:proj} with $\mat{W} \in \R^{256 \times 1024}$: at each
position it learns how much of each scale to use. Splitting $\mat{W}$ into four $256 \times 256$ blocks
$[\mat{W}^{(1)}\ \mat{W}^{(2)}\ \mat{W}^{(3)}\ \mat{W}^{(4)}]$ shows what it does:
\[
\mat{W} \begin{pmatrix} \vect{p}^{(1)}\\ \vect{p}^{(2)}\\ \vect{p}^{(3)}\\ \vect{p}^{(4)} \end{pmatrix}
= \mat{W}^{(1)}\vect{p}^{(1)} + \mat{W}^{(2)}\vect{p}^{(2)} + \mat{W}^{(3)}\vect{p}^{(3)} + \mat{W}^{(4)}\vect{p}^{(4)} :
\]
a learned sum of the four scales, one linear map per scale. The nonlinearity $\GELU$
(\cref{ch:block}) follows. The result: one 256-channel map at 32 $\times$ 32 carrying all four scales.
This projection, upsample, concatenate and fuse pattern is what the SegFormer paper calls an
\index{all-MLP decoder}\emph{all-MLP decoder}, because 1 $\times$ 1 convolutions are per-pixel MLP layers.

\begin{keypoint}
Up to the fuse, the decoder is a per-pixel linear mix of four interpolated maps; everything spatial in
it so far is the smoothing of bilinear interpolation.
\end{keypoint}

\section{Step 5: refine at quarter resolution, then pixel shuffle}\label{sec:decoder:shuffle}

\begin{lstlisting}
self.refine = nn.Sequential(
    nn.Conv2d(decoder_dim, decoder_dim, kernel_size=3, padding=1),          # 256 -> 256
    nn.GELU(),
    nn.Conv2d(decoder_dim, out_channels * 4 * 4, kernel_size=3, padding=1)) # 256 -> 512
...
out = F.pixel_shuffle(refined, 4)          # (B, 512, 32, 32) -> (B, 32, 128, 128)
\end{lstlisting}

\subsection{The two 3 \texorpdfstring{$\times$}{x} 3 convolutions}

Both keep the grid at 32 $\times$ 32: by the output-size formula of \cref{ch:encoder} with $k = 3$, $s =
1$, $p = 1$, a side of 32 becomes $\lfloor (32 + 2 - 3)/1 \rfloor + 1 = 31 + 1 = 32$. They are the
decoder's only spatial layers. A 3 $\times$ 3 window at quarter resolution spans 3 cells of 4 pixels,
$3 \times 4 = 12$ original pixels across; two in a row reach $1 + 2 \cdot 2 = 5$ cells, 20 pixels (the
receptive-field recurrence of \cref{ch:encoder} with $j = 4$). That is local context around every
point at one sixteenth of the cost of full resolution.

The second convolution, a \index{sub-pixel convolution}\emph{sub-pixel convolution}, outputs $D \times 4 \times 4 = 32 \times 16 = 512$ channels, where $D = 32$ is the
compressed band count, and has \textbf{no activation after it}. Its output flows through the
decompressor into denormalisation and must be free to take any value, including large negatives; a
GELU or sigmoid would make whole ranges of reflectance unreachable. The same reasoning keeps the
decompressor bare (\cref{ch:compressor}).

\subsection{Pixel shuffle: the index map}

\begin{definition}[Pixel shuffle]
For an upscale factor $r$, \code{pixel_shuffle} maps a tensor of shape $(B, C r^2, H, W)$ to $(B, C, rH,
rW)$ by
\begin{equation}\label{eq:decoder:shuffle}
\text{out}[b, c,\ rh + i,\ rw + j] = \text{in}[b,\ c r^2 + r i + j,\ h, w],
\end{equation}
for $0 \le c < C$, $0 \le h < H$, $0 \le w < W$ and $0 \le i, j < r$.
\end{definition}

Read it from right to left. Input channels come in $C$ consecutive groups of $r^2$; group $c$ becomes
output channel $c$. Inside a group, channel offset $ri + j$ is a row-major position $(i, j)$ in an $r
\times r$ block (row $i = \lfloor \text{offset}/r \rfloor$, column $j = \text{offset} \bmod r$). That
block is placed at rows $rh .. rh + r - 1$ and columns $rw .. rw + r - 1$, the output block that sits
over input position $(h, w)$.

\paragraph{Inverting it.} Given an output pixel $(c, Y, X)$, which input number made it? Divide each
coordinate by $r$ with remainder: $h = \lfloor Y/r \rfloor$, $i = Y \bmod r$, $w = \lfloor X/r \rfloor$,
$j = X \bmod r$. Then the source is channel $c r^2 + r i + j$ at position $(h, w)$. Every output pixel
has exactly one source and every input number has exactly one destination: pixel shuffle is a
permutation of the numbers, with no arithmetic and no parameters.

\begin{worked}[label={we:decoder:shuffle}]{Pixel shuffle by hand, and PyTorch's order}
Take $C = 2$, $r = 2$, $H = W = 1$: an input of $C r^2 = 8$ channels, each a single number. Let channel
$k$ hold the value $k$, so in $= [0, 1, 2, 3, 4, 5, 6, 7]$. Apply \eqref{eq:decoder:shuffle} with $h = w
= 0$, so the output rows are $i$ and columns $j$. On paper:
\begin{align*}
c = 0:&\quad (i,j) = (0,0) \leftarrow 0\cdot 4 + 0 + 0 = 0, & (0,1) &\leftarrow 0 + 0 + 1 = 1,\\
      &\quad (1,0) \leftarrow 0 + 2 + 0 = 2, & (1,1) &\leftarrow 0 + 2 + 1 = 3,\\
c = 1:&\quad (0,0) \leftarrow 1\cdot 4 + 0 + 0 = 4, & (0,1) &\leftarrow 4 + 0 + 1 = 5,\\
      &\quad (1,0) \leftarrow 4 + 2 + 0 = 6, & (1,1) &\leftarrow 4 + 2 + 1 = 7.
\end{align*}
\[
\text{out}[0] = \begin{pmatrix} 0 & 1\\ 2 & 3\end{pmatrix}, \qquad
\text{out}[1] = \begin{pmatrix} 4 & 5\\ 6 & 7\end{pmatrix}.
\]
Running \code{F.pixel_shuffle} on this input in PyTorch 2.5.1 gives exactly these two blocks
(\file{research/textbook/tools/compute_masking_output.txt}), which confirms the ordering: channels are
grouped by output channel first, and within a group the offset runs along a row before moving down.

Now the real case, $r = 4$, $C = 32$. Which input number produces output pixel $(c, Y, X) = (1, 5, 6)$?
\begin{align*}
h &= \lfloor 5/4 \rfloor = 1, & i &= 5 \bmod 4 = 1,\\
w &= \lfloor 6/4 \rfloor = 1, & j &= 6 \bmod 4 = 2,\\
\text{channel} &= 1 \cdot 16 + 4 \cdot 1 + 2 = 22 .
\end{align*}
So the source is input channel 22 at position $(h, w) = (1, 1)$.
\end{worked}

The count of numbers is unchanged, as it must be for a permutation:
\begin{align*}
\text{before:}\quad 512 \times 32 \times 32 &= 512 \times 1024 = 524{,}288,\\
\text{after:}\quad 32 \times 128 \times 128 &= 32 \times 16{,}384 = 524{,}288 .
\end{align*}

\coursefig[0.95\linewidth]{iw-21-decoder}{The decoder. The four encoder maps are projected to 256
channels by 1 $\times$ 1 convolutions, resized to F1's 32 $\times$ 32, concatenated to 1,024 channels,
fused and refined, and expanded to $512 = 32 \times 4 \times 4$ channels. PixelShuffle(4) rearranges
each group of 16 channels into a 4 $\times$ 4 block of pixels (inset, with $r = 2$), returning
\shape{(B, 32, 128, 128)} without any arithmetic.}{fig:decoder:course}

\subsection{Why pixel shuffle and not a final bilinear upsampling}

With pixel shuffle, output pixel $(i, j)$ of a 4 $\times$ 4 block comes from its own channel $16c + 4i
+ j$, and that channel is produced by its own row of weights in the last convolution. So the 16 pixels
of a block can take 16 unrelated values. With bilinear upsampling all 16 would be weighted averages of
the same few coarse values.

\begin{worked}[label={we:decoder:spike}]{A one-pixel anomaly, two ways}
One row of a 4-pixel block, true values $[0, 0, 8, 0]$ (one hot pixel on a cold background, in
illustrative units). A bilinear upsampler at factor 4 sees one coarse value per block, the block mean
\[
\tfrac14 (0 + 0 + 8 + 0) = 2,
\]
and neighbouring blocks with value 0. Whatever it does, it outputs a smooth bump between neighbouring
coarse values: every pixel of the block gets a value between 0 and 2. The prediction at the hot pixel is
at most 2, the residual there at least $8 - 2 = 6$, and the three cold pixels pick up spurious residuals
of up to 2 each. A pixel-shuffle head has four channels for this row, one per pixel, and can output
$[0, 0, 8, 0]$ exactly if its weights have learned that pattern.
\end{worked}

The decoder's own comment puts it in the terms of the thermal model it was written for: ``a 1-pixel
anomaly at 45\textdegree C surrounded by 30\textdegree C background gets its own predicted value, not a
smoothed average of the 4x4 block'' (\file{seg_former_decoder.py}, lines 107--109). The important
subtlety: you do \emph{not} want the model to reconstruct an anomaly well. You want it to be
\emph{able} to give each pixel its own prediction, so that what it predicts at an anomalous pixel is
``what normally goes here'', unblurred, and the difference is large and sharp.

\subsection{What it costs}

A convolution's work\index{multiply-add} is its multiply-adds per position (the weight count without biases) times the
number of positions. For the two refine convolutions at 32 $\times$ 32:
\begin{align*}
\text{refine 1:}\quad 256 \times 256 \times 9 &= 589{,}824 \text{ per position}, & 589{,}824 \times 1024 &= 603{,}979{,}776,\\
\text{refine 2:}\quad 256 \times 512 \times 9 &= 1{,}179{,}648 \text{ per position}, & 1{,}179{,}648 \times 1024 &= 1{,}207{,}959{,}552,
\end{align*}
about 1.8 billion multiply-adds per patch\tagC. At 128 $\times$ 128 the same two convolutions would do
sixteen times as much, about 29 billion. Pixel shuffle itself costs nothing but memory movement.

\begin{sourcenote}[: thermal comments in a hyperspectral decoder]
\code{SegFormerDecoder} was written for the thermal model (Chakshu) and is reused unchanged: its default
is \code{out_channels=1} and its comments speak of temperatures (lines 49--51, 107--112). The
hyperspectral model builds it with \code{out_channels=compressed_channels}
(\file{hyperspectral_seg_former_mae.py}, lines 107--111), 32 in every shipped configuration, so the last
convolution has $32 \times 16 = 512$ outputs. The same class, built with width 128 and 5 output
channels, is the segmentation head of \cref{ch:seghead}; a different, progressive decoder is the subject
of \cref{ch:decoderwork}.
\end{sourcenote}

\section{Decompress and denormalise}\label{sec:decoder:back}

The decoder's output is still in the compressed 32-channel space. Two last steps bring it back to
bands and to physical units.

\paragraph{The decompressor}\index{decompressor} (\cref{ch:compressor}) is a 1 $\times$ 1 convolution from 32 to 165
channels, with no normalisation and no activation:
\[
\shape{(B, 32, 128, 128)} \xrightarrow{\text{Conv2d(32, 165, k=1)}} \shape{(B, 165, 128, 128)} .
\]
At each pixel it computes $\hat{\vect{z}} = \mat{W}_{\text{dec}}\, \vect{y} + \vect{b}_{\text{dec}}$ with
$\mat{W}_{\text{dec}} \in \R^{165 \times 32}$: each of the 165 bands is a learned mixture of the 32
compressed channels.

\paragraph{Denormalisation}\index{denormalisation} undoes the per-band z-score of \cref{ch:bandpipeline}. With the band's
stored mean $\mu_c$ and standard deviation $\sigma_c$ (buffers baked into the model):
\begin{equation}\label{eq:decoder:denorm}
\hat{x}_{c,y,x} = \sigma_c\, \hat{z}_{c,y,x} + \mu_c .
\end{equation}
It is the normalisation $z = (x - \mu)/\sigma$ solved for $x$: multiply both sides by $\sigma$, then add
$\mu$.

\begin{worked}[label={we:decoder:denorm}]{Denormalising one value}
Suppose a band has $\mu = 0.10$ and $\sigma = 0.08$\tagA{} (illustrative values of the right size for
reflectance), and the decompressor outputs $\hat{z} = -0.5$ at some pixel. Then
\begin{align*}
\hat{x} &= 0.08 \times (-0.5) + 0.10 && \text{multiply by } \sigma, \text{ then add } \mu\\
        &= -0.04 + 0.10\\
        &= 0.06 .
\end{align*}
A predicted reflectance of 0.06, half a standard deviation below the band mean.
\end{worked}

\section{One patch, all the way through}\label{sec:decoder:trace}

Every component of the model's \code{forward} method is now yours
(\file{app/foundation_models/components/hyperspectral_seg_former_mae.py}, lines 116--156). \index{forward pass}\index{shape-preserving model}Here is the whole walk for a batch of $B = 4$ patches of 128 $\times$ 128 with 165 bands,
compressed to $D = 32$, at $\rho = 0.65$.

\begin{lstlisting}
def forward(self, x, mask=None, keep_mask=None):   # x (4,165,128,128), mask (4,1,128,128), keep_mask (4,1024)
    if mask is not None:
        x = x * mask                                # invalid pixels -> 0.0 in every band
    if self.normalize is not None:
        x = self.normalize(x)                       # per-band z-score
    x = self.compressor(x)                          # (4, 32, 128, 128)
    features = self.encoder(x, keep_mask=keep_mask) # [F1, F2, F3, F4]
    x_hat = self.decoder(features)                  # (4, 32, 128, 128)
    x_hat = self.decompressor(x_hat)                # (4, 165, 128, 128)
    if self.denormalize is not None:
        x_hat = self.denormalize(x_hat)             # back to reflectance
    return x_hat
\end{lstlisting}

Notice what is \emph{not} here: no loss, no score, no masking strategy. The model reconstructs whatever
it is given with whatever keep mask it is handed. The trainer chooses random masks (\cref{ch:masking}),
the inferencer chooses complementary ones (\cref{ch:inference}); the same class serves both.

\begin{longtable}{@{}rlll@{}}
\caption{One forward pass, every shape. $B = 4$, 165 bands, $D = 32$, 128 $\times$ 128, $\rho = 0.65$
on fully valid patches. Shapes confirmed by running the real model on CPU
(\texttt{research/textbook/tools/compute\_params.py}).}\label{tab:decoder:trace}\\
\toprule
\# & Operation & Output shape & Note \\\midrule
\endfirsthead
\toprule
\# & Operation & Output shape & Note \\\midrule
\endhead
& input $x$ & \shape{(4, 165, 128, 128)} & reflectance \\
1 & \texttt{x * mask} & \shape{(4, 165, 128, 128)} & mask broadcasts over bands \\
2 & normalise & \shape{(4, 165, 128, 128)} & z-scores; invalid $\approx -\mu/\sigma$ \\
3 & compressor & \shape{(4, 32, 128, 128)} & 1 $\times$ 1 conv + BatchNorm \\\midrule
4 & S1 patch embed & \shape{(4, 1024, 32)} & $k4\,s4\,p0$, grid 32 $\times$ 32 \\
  & S1 remove tokens & \shape{(4, 359, 32)} & 665 hidden per patch \\
  & S1 2 blocks + LN & \shape{(4, 359, 32)} & full-attention fallback \\
  & S1 restore + reshape & \shape{(4, 32, 32, 32)} & F1, zeros at hidden tokens \\
  & S2 embed, blocks, reshape & \shape{(4, 64, 16, 16)} & F2, $R = 4$ \\
  & S3 embed, blocks, reshape & \shape{(4, 160, 8, 8)} & F3, $R = 2$ \\
  & S4 embed, blocks, reshape & \shape{(4, 256, 4, 4)} & F4, $R = 1$ \\\midrule
5 & 4 projections & \shape{(4, 256, H\_i, W\_i)} & one 1 $\times$ 1 conv per map \\
  & bilinear to 32 $\times$ 32 & $4 \times$ \shape{(4, 256, 32, 32)} & F1 skipped \\
  & concatenate & \shape{(4, 1024, 32, 32)} & $4 \times 256$ \\
  & fuse + GELU & \shape{(4, 256, 32, 32)} & 1 $\times$ 1 conv \\
  & refine 1 + GELU & \shape{(4, 256, 32, 32)} & 3 $\times$ 3 conv \\
  & refine 2 & \shape{(4, 512, 32, 32)} & $512 = 32 \times 16$, no activation \\
  & pixel shuffle(4) & \shape{(4, 32, 128, 128)} & a permutation \\\midrule
6 & decompressor & \shape{(4, 165, 128, 128)} & 1 $\times$ 1 conv, bare \\
7 & denormalise & \shape{(4, 165, 128, 128)} & $\sigma_c \hat z + \mu_c$ \\\bottomrule
\end{longtable}

\coursefig[0.98\linewidth]{iw-22-forward}{The seven steps of \texttt{forward()} for a batch of four 128
$\times$ 128 patches. Spectral steps are teal, spatial steps indigo. Step 4 expands into the four encoder
stages (hiding at stage 1 only). The model stops at the reconstruction; the trainer computes the loss
(amber) in raw reflectance.}{fig:decoder:forward}

The shape story on one line:
\[
(4, 165, 128, 128) \to (4, 32, 128, 128) \to \text{tokens} \to \text{F1..F4} \to (4, 32, 128, 128) \to (4, 165, 128, 128).
\]
Input and output shapes are identical. Reconstruction models are shape-preserving; they have to be,
or you could not subtract one from the other.

\subsection{What the caller does next, and in which units}

The trainer takes the output and builds the loss mask of \cref{ch:masking}:
\begin{lstlisting}
x_hat = model(pixels, mask=mask, keep_mask=keep_mask)          # (B, 165, H, W), reflectance
loss_mask = pixel_pred_mask * eroded_mask                       # (B, 1, H, W)
per_pixel_l1_mean = (x_hat - pixels).abs().mean(dim=1, keepdim=True)   # (B, 1, H, W)
l1_loss = per_pixel_l1_mean[loss_mask == 1].mean()
\end{lstlisting}
It compares the \emph{denormalised} output with the \emph{raw} input: reflectance against reflectance.
So an L1 of 0.008 means that, averaged over graded pixels, each band's predicted reflectance is off by
0.008, a few percent of typical values between 0.02 and 0.30. In normalised space the same number would
mean ``0.008 standard deviations'' and would need translating before it meant anything. The losses
themselves are \cref{ch:losses}.

\begin{drift}
The L1 of 0.008 is the upper end of the expectation ``\textasciitilde 0.003--0.008'' at epoch 500 written into the
v0.2.0 configuration's notes\tagM{configs/hyperspectral_segformer_exp_2.json} (line 15), not a measured
result of the shipped model; the same file states that the loss is computed in raw reflectance space
(line 6). Measured losses for each run are in \cref{ch:pretrain} and \cref{ch:ftruns}.
\end{drift}

\begin{confusions}
\confusion{Why upsample to F1's size rather than to full resolution?}{Cost. F1's grid is a quarter of the
patch side, so convolutions there are sixteen times cheaper. The final jump is a free rearrangement.}
\confusion{Is pixel shuffle learned?}{No. It has no parameters; it moves numbers. What is learned is the
convolution before it, which produces the 16 sub-pixel channels per output channel.}
\confusion{Does the decoder know which positions were hidden?}{Not explicitly. F1 arrives with zero
vectors there, an implicit signal; no mask is passed to the decoder.}
\confusion{Why is it fine that F1 has holes?}{F2, F3 and F4 cover the same ground without holes, and the
fuse convolution learns how much to take from each scale at each position.}
\confusion{Then why use bilinear interpolation at all?}{For the coarse maps it is the right tool: their
content is smooth context, and smoothing it onto F1's grid loses nothing. Only the last jump to
individual pixels must not smooth.}
\confusion{Compress, then decompress: does it undo itself?}{No. The encoder and decoder sit between them
and change the representation. The compressor makes the transformer affordable; the decompressor turns
its answer back into bands.}
\end{confusions}

\begin{checkyourself}
\item\label{ex:decoder:1} Upsample the row $[2, 10]$ to four values with \code{align_corners=False},
showing $u_d$ for each output. Then upsample the 2 $\times$ 2 grid $\begin{pmatrix} 0 & 4\\ 8 &
12\end{pmatrix}$ to 4 $\times$ 4.
\item\label{ex:decoder:2} Apply \code{pixel_shuffle} with $r = 2$ to a 1 $\times$ 1 input with 4
channels holding $[9, 8, 7, 6]$.
\item\label{ex:decoder:3} With $r = 4$ and $C = 32$, which input channel and position produce output
pixel $(c, Y, X) = (31, 127, 0)$? And which output pixel does input channel 37 at position $(2, 3)$ go
to?
\item\label{ex:decoder:4} If the compressed width were $D = 24$, what would the last refine convolution's
output channel count be, and how many numbers would the tensor before pixel shuffle hold for one
patch? Check it against the decoder's output size.
\item\label{ex:decoder:5} Give the shape after each of the seven steps of \code{forward} for a batch of
2 patches of 256 $\times$ 256 with 165 bands and $D = 32$, including F1--F4.
\item\label{ex:decoder:6} A band has $\mu = 0.25$, $\sigma = 0.1$. The decompressor outputs $\hat z =
1.3$. What reflectance does the model predict? If the true reflectance is 0.36, what is that pixel's
absolute error in this band?
\end{checkyourself}

\nextchapter{ch:params}
